% !PRINTED 606
\weberpage{606}
%<<CONT
% [page alignment mirrored 2026-08-07] Printed 606 begins with this
% continuation from i0630 and ends after the §179 footnote.  The former
% "We assume ..." tail is printed on 607 and already belongs to i0632.
then $f(x)$ is irreducible, and one still has freedom enough in the choice of the integers $c$ to bring the rational fractions $a_1, a_2 \ldots a_n$ arbitrarily near to any arbitrarily given system of values.

But with this the theorem is proved:

4. \gesperrt{Of every arbitrary prime degree there are infinitely many equations with rational coefficients without affect.}

The proof that has been given here shows that such affect-free equations can be found whose coefficients fill a finite region everywhere densely — namely the region in which the real coefficients of the equations with only two
conjugate imaginary roots lie. The proof is therefore, apart from the fact that it is restricted to prime degrees, also not exhaustive in this respect, that it gives us no information about the remaining regions, in which in all probability the matter stands just the same.

\begin{center}
\S.\,179.

\gesperrt{Solution by real radicals.}
\end{center}

In the solution of the cubic equations with real coefficients, the two cases in which the discriminant is negative or positive have been distinguished from ancient times. In the second case, although the equation then has precisely three real roots, one has, in the application of the \gesperrt{Cardan formula}, to extract the cube root of an imaginary expression, and the attempt to represent the roots in real form leads always back to a cubic equation of the same character. For this reason this case has been called the \gesperrt{casus irreducibilis}. The cubic equations of the casus irreducibilis are a special case of the \gesperrt{cyclic equations with real roots}, which we came to know in \S.\,165, and in whose solution likewise roots of imaginary quantities, or what comes to the same thing, divisions of angles, occur. That these roots of imaginary quantities or divisions of angles can in no way be avoided, we can now prove.\footnote{Cf. \gesperrt{H\"older}, Mathematische Annalen, Bd. 38; \gesperrt{Kneser}, ebendas., Bd. 41.}
